\clearpage
\fancyhead[L]{\small CONNAISSANCE DES TEMS.}
\hypertarget{native-559-lower}{}\label{native-559-lower}
\begin{center}\small Göttingische gelehrte Anzeigen. Stück 8. Seite 873..876. 1813 Juni 3.\\[3pt]\rule{.72\linewidth}{.3pt}\end{center}

\emph{Connaissance des tems ou des mouvemens célestes, à l'usage des astronomes et des navigateurs, pour l'an} 1813, publiée par le bureau des longitudes. Juillet 1811. Paris. 228 Seiten in Octav.

\emph{Connaissance des tems etc.} pour l'an 1814. Avril 1812. 240 Seiten.

Die Tafel für die geographische Lage der vornehmsten Oerter der Erde hat besonders durch die Benutzung der auf den Reisen von \textsc{Dentrecasteaux}, von \textsc{Humboldt} und von \textsc{Krusenstern} gemachten Bestimmungen an Umfang wieder beträchtlich gewonnen, so dass sie in dem Jahrgange für 1814 über 1600 Artikel, und folglich über 100 mehr enthält, als in dem Jahrgange für 1812. Die \emph{Zusätze} zu diesen beiden Jahrgängen sind diesmal nicht zahlreich. Der Jahrgang 1813 enthält blos zwei kleine Aufsätze von \textsc{Laplace}. In dem ersteren gibt er eine wichtige Bereicherung der Theorie der so genannten Methode der kleinsten Quadrate, indem er durch eine künstliche Analyse zeigt, dass bei derjenigen Combination der Grundgleichungen, welche diese Methode vorschreibt, der zu befürchtende mittlere Fehler in den Resultaten ein Minimum wird, so bald die Anzahl der Grundgleichungen sehr gross ist, das Gesetz der Wahrscheinlichkeit der Fehler sei, welches es wolle. Umständlicher ist diese Untersuchung in dem unlängst erschienenen grössern Werke des Verf. über die Wahrscheinlichkeitsrechnung ausgeführt. Der andere Zusatz betrifft die Mondsgleichung von langer Periode, deren Nothwendigkeit man aus der Disharmonie zwischen den für 1692, 1756, 1779 und 1801 herausgebrachten Epochen der mittleren Mondslänge \SourceLeaf{560}geschlossen hat, ohne sie bisher ganz befriedigend in der Theorie nachweisen zu können. In den \textsc{Bürg}schen Mondstafeln war sie dem Producte des Sinus der doppelten Länge des Mondsknoten plus der Länge des Mondsperigeum weniger der dreifachen Länge des Sonnenperigeum proportional, und ihr Coëfficient $= -14''$ angenommen worden. \textsc{Laplace} kündigt jetzt an, dass er überwiegende Gründe habe (über welche er sich aber nicht näher erklärt), sie dem Cosinus der doppelten Länge des Mondsknoten plus der Länge des Mondsperigeum proportional zu setzen, und bestimmt ihren Coëfficienten $= -13''96$. (\textsc{Burckhardt} hat in seinen neuen Mondstafeln, wie wir bereits bei deren Anzeige bemerkten. dieselbe Form zum Grunde gelegt und den Coëfficienten $= -12''5$ angenommen.) Wir sehen mit Verlangen der ausführlicheren Darstellung der theoretischen Untersuchungen des Verf. über diese delicate Frage entgegen.

Die Zusätze zu dem Jahrgange 1814 sind folgende: 1) Ueber den Ursprung der Cometen, von \textsc{Lagrange}. Der grosse (den Wissenschaften seitdem, wenn gleich in hohem Alter, doch noch immer zu früh, durch den Tod entrissene) Geometer dehnt hier die bekannte \textsc{Olbers}'sche Hypothese über den Ursprung der neuen Planeten auf die Cometen aus, und entwickelt mit derjenigen Eleganz, die man immer von ihm gewohnt war, die Bedingungen, unter denen ein von einem Planeten, der sich in einer Kreisbahn bewegte, fortgeschleudertes Fragment eine parabolische Bahn beschreibt. Die Grenzen der relativen Geschwindigkeit, womit der werdende Comet fortgestossen wird, sind, wie man leicht sich überzeugt (\textsc{Lagrange} hat es indess nicht bestimmt ausgesprochen), $\BareRoot{2}-1$ und $\BareRoot{2}+1$, die Geschwindigkeit des Planeten als Einheit angesehen; merkwürdig ist aber der hier von \textsc{Lagrange} aufgestellte (gleichfalls leicht zu beweisende) Satz, dass die Geschwindigkeit $\BareRoot{3}$ die Scheidewand zwischen den rechtläufigen und rückläufigen Cometen bildet, so dass ein (in Beziehung auf die Bahn des Planeten) rechtläufiger eine kleinere, der rückläufige eine grössere erfordert. Denkt man sich einen Planeten in einer hundert Mal so grossen Entfernung, wie unsere Erde, der durch die Wirkung eines in seinem Innern plötzlich, durch was immer für Ursachen, frei werdenden elastischen Fluidum in mehrere Stücke zersprengt würde, so brauchte die Explosion nur so stark zu sein, dass sie eine zwölf oder funfzehn Mal so grosse Geschwindigkeit, wie die einer 24pfündigen Kanonenkugel, ertheilen könnte, um elliptische oder parabolische Cometen nach allen möglichen Dimensionen und nach \emph{allen möglichen Richtungen hervorzubringen}. Dieser letztern Behauptung müssen wir indess widersprechen; es könnten auf diese Art nur solche Cometen entstehen, deren Knotenlinie auf der vorigen Planetenbahn mit der Apsidenlinie sehr nahe zusammenfällt, wenn sie in ihrer Sonnennähe bis in die Region der Erdbahn herabkommen sollen. Allein von den sämmtlichen bisher berechneten Cometenbahnen ist nur ein kleiner Theil von der Art, dass die Möglichkeit einer solchen Entstehungsart zugestanden werden könnte. Praktische Astronomen werden überhaupt schwerlich geneigt sein, in Cometen, deren ganzes Ansehen auf eine durchaus verschiedene physische Beschaffenheit hindeutet, Stücke von festen Planetenkörpern zu erkennen. 2) Einige Notizen aus \textsc{Krusenstern}'s Reise um die Welt. 3) Eine sehr bequem angeordnete Tafel zur Verwandlung der Sternzeit in mittlere Sonnenzeit, von \textsc{Burckhardt}. 4) Auszug aus den meteorologischen Beobachtungen zu Paris im Jahre 1809. Am Schlusse beider Jahrgänge, wie gewöhnlich, das Verzeichniss der Mitglieder des Längen-Büreau.
